% !TeX root = 20261029_exp91.tex
% !LW recipe=lablog XeLaTeX
\input{_support/style.tex}
\begin{document}
\reporttitle{2026-10-29}{EXP91: 두 TabPFN 설계의 실제 IDS 분류 결과와 실행 기록}
\subsection{완료 결과와 보고서 복구}
두 설계의 실제 Intrusion Detection System (IDS) 분류 실험은 10개 조건, 총 20개 모델 평가를 완료했다.
실행은 2026-10-29 01:18:16부터 04:36:46까지 한국 시각으로 진행됐으며 연구 버전은 v7, 실험은 EXP91이다.
원격 장비에서는 보고서 원문 작성 후 \texttt{kotex.sty} 누락으로 PDF 컴파일이 실패했다.
이 장치에서 원문과 결과를 확인하고, 저장된 혼동행렬로 지표를 재계산한 뒤 임시 PDF 컴파일을 검증한다.
이 앞부분은 완료 결과의 해석이며, 뒤의 준비·실패 기록은 당시 상태를 보존한 부록이다.

\sourcepath{tabpfn/results/v7/exp91/20261028_161815_remote_both_s43_fixedenv_ecd34e66/summary.json}

\subsection{비교 대상과 평가 설정}
1안(A)은 고정 TabPFN 교사의 확률과 비라벨 특징을 이용해 학생 인코더·대표점을 학습하고, 최종 분류에는 학생만 사용한다.
2안(B)은 전체 라벨의 Global TabPFN에 혼동 클래스 쌍 전문가를 추가하며, Out-of-Fold (OOF) 예측에서 선택 조건을 충족한 보정만 채택한다.
두 안의 정상 라벨은 750개로 고정하고 공격 클래스별 라벨 예산은 $n\in\{2,3,5,10,20\}$, seed는 43이다.
A는 추가로 학습 분할의 비라벨 10,000개를 사용하지만 B는 이를 사용하지 않으므로, 두 안의 정보 예산 차이를 결과 해석에 포함한다.
TabPFN의 가중치는 고정하며 설정한 앙상블 수는 4, 추론 배치 크기는 1,024이다.

정제 전 NF-v3의 정상 포함 전체 test는 CIC2018 4,023,114행, ToN-IoT 5,504,060행이다.
공격 Macro-F1은 정상 클래스의 F1을 제외한 클래스별 평균이지만, 정상 행을 포함한 전체 혼동행렬에서 계산하므로 정상의 공격 오분류가 공격 정밀도에 반영된다.
False Positive Rate (FPR)는 실제 정상 중 공격으로 분류된 비율이다.
이 보고서의 재검산은 저장된 혼동행렬의 클래스별 표본 수와 지표를 확인하며, Git에서 제외된 행별 예측을 다시 생성하거나 원본 정답과 재대조한 검증은 아니다.

\subsection{공격 분류와 정상 오탐 결과}
1안은 10개 조건 중 6개에서 2안보다 높은 공격 Macro-F1을 기록했지만, 정상 오탐률도 10개 조건 모두에서 높았다.
예를 들어 CIC2018 $n=10$에서 A/B의 공격 Macro-F1은 65.12/49.68\%, 정상 오탐률은 1.165/0.079\%였다.
ToN-IoT $n=20$에서는 A/B의 공격 Macro-F1이 48.98/49.24\%, 정상 오탐률이 9.917/8.254\%였다.
따라서 공격 F1만으로 1안의 전체 탐지 이점이나 정상 오탐을 고려한 우위를 확정하지 않는다.
표의 F1과 FPR은 모두 백분율이며, 전문가 수는 최종 B에 실제 포함된 이진 모델 수다.

\input{evidence/20261029/exp91_results_table.tex}
\reportfigure{figures/20261029/exp91_remote_results.pdf}{정제 전·정상 750개 고정·seed 43의 실제 전체 test 결과. 위는 공격 Macro-F1, 아래는 정상 오탐률이다. 빨간 실선은 1안(A), 파란 점선은 2안(B)이며 단일 seed이므로 오차막대는 제시하지 않는다.}

\subsection{전문가 채택과 비용 해석}
2안의 쌍 전문가는 3개 조건에서만 채택됐으며, 나머지 7개 조건은 Global TabPFN만으로 최종 예측을 수행했다.
CIC2018 $n=10$에서는 정상--DoS 한 쌍, ToN-IoT $n=10$에서는 injection--scanning 한 쌍, ToN-IoT $n=20$에서는 backdoor--scanning과 injection--ransomware 두 쌍이 채택됐다.
$n=2$에서는 쌍 선택에 필요한 클래스별 최소 3개 독립 특징 그룹 조건을 만족하지 못했다.
전문가를 채택한 3개 조건에는 같은 입력의 Global-only 전체 test 결과가 별도로 저장되지 않았으므로, 현재 자료만으로 전문가 추가의 순수 증분 효과를 계산하지 않는다.
OOF 선택 지표의 개선을 전체 test의 개선으로 대신 해석하지 않는다.

1안의 전체 test 추론 시간은 5.06--10.73초, 2안은 596.75--2,517.01초로 기록됐다.
1안의 fit 시간 31.82--71.10초에는 교사 확률 생성과 학생 학습이 포함되며, 빠른 최종 추론이 이 준비 비용을 없애는 것은 아니다.
이 시간들은 NVIDIA A100-SXM4-40GB에서 기록한 모델 fit·추론 구간이며, 입력 로딩·모델 저장·복원 확인 시간은 해당 구간 밖이다.
초기에 중복 실행이 같은 GPU에서 겹쳤으므로 시간 값은 해당 실행의 관측치로만 사용하고, 통제된 단독 실행의 속도 배율로 주장하지 않는다.

\subsection{현재 근거와 후속 개선}
정상 오탐을 함께 줄이면서 희귀 공격을 보완하는지가 후속 설계 판단의 기준이다.
1안은 정상 오탐 제약을 포함한 학습 목적이나 정상--공격 분리 단계를 후보로 두고, 교사만 사용·학생 감독학습만 사용·비라벨 구성 제거 절제를 같은 라벨 예산에서 비교한다.
2안은 전문가를 채택한 조건의 Global-only 전체 test 평가를 우선 추가하고, 극소수 독립 그룹에서의 선택 불확실성을 반영하는 축소 보정이나 보수적 전문가 공유를 검토한다.
이 개선안의 효과는 아직 검증되지 않았으며, 현재 test에 맞춰 임계값을 정한 뒤 같은 test를 확증 근거로 재사용하지 않는다.

단일 seed 결과와 서로 다른 비라벨 접근만으로 제안 구조의 일반적 우위를 확정할 수 없다.
이번 실행에는 XGBoost나 문헌의 경쟁 모델 비교가 포함되지 않았고, 과거 검토한 분할의 학습/비라벨--test 동일 특징 중복도 유지됐다.
후속 비교는 동일한 정보 예산의 대조와 추가 seed, 별도로 남긴 평가 분할에서 수행하고, 현재 결과는 탐색적 설계 진단으로 보관한다.
완료 표식과 지표 재검산은 실행의 성공을 확인하지만 논문 주장 채택을 뜻하지 않는다.

\sourcepath{docs/lablog/evidence/20261029/exp91_completed_run_validation.json}
\sourcepath{docs/lablog/evidence/20261029/exp91_results.csv}
\sourcepath{docs/lablog/evidence/20261029/exp91_class_f1.csv}

\clearpage
\section*{부록: 준비 과정과 당시 실행 기록}
아래 기록의 성능 미실행·재실행 필요 표기는 각 준비·실패 단계 당시의 상태다.
최종 완료 상태와 현재 사용할 결과는 앞부분의 실제 분류 결과를 기준으로 한다.
\subsection{목적과 수행 범위}
원격 장비에서 입력 경로를 추측하지 않고 두 설계의 분류 실험을 실행할 수 있도록 입력 생성기와 야간 실행기를 연결한다. 사용자는 해당 장비에 원본 data와 checkpoint가 이미 있다고 확인했다. 이전 안내의 실제\_\allowbreak{}실험입력.npz는 존재하는 파일이 아니었으며, 이번 작업에서 실제 입력 생성 단계를 보완했다. 이 기록은 2026-10-29에 수행한 자료 준비·코드 검증이며, 모델 분류 성능을 측정한 결과가 아니다.

Intrusion Detection System (IDS; 침입 탐지 시스템) 자료는 기존 정제 전 NF-v3 suite를 그대로 사용한다. source는 data/\allowbreak{}nfv3\_\allowbreak{}energy\_\allowbreak{}suite\_\allowbreak{}uncapped\_\allowbreak{}scenarios.pkl이며 내용 해시를 EXP89 기준으로 검증한다. data 폴더·NPZ·checkpoint는 Git 제외 대상이고, 두 실행 스크립트와 설정만 Git으로 배포한다. 원격 장비에서는 기존 원본을 이용해 NPZ를 재생성한다.

\subsection{입력 위치와 샘플 구성}
입력은 data/\allowbreak{}exp91/\allowbreak{}inputs/\allowbreak{}에 생성한다. 파일명은 데이터셋\_\allowbreak{}original\_\allowbreak{}fixed750\_\allowbreak{}n라벨수\_\allowbreak{}s시드\_\allowbreak{}u비라벨수\_\allowbreak{}full.npz 형식이며 같은 이름의 JSON에 클래스 이름·정상 ID·범주 열·행 수·해시를 저장한다. 기본 명령은 CIC2018·ToN-IoT, 공격별 n=\allowbreak{}2·3·5·10·20, seed 43의 10개 조건을 만든다. seed 43--47 확장은 두 스크립트 모두 --seeds 43 44 45 46 47로 지정한다.

학습 라벨은 EXP89의 클래스별 20개 추출과 중첩 prefix를 재현하고 정상은 EXP90의 고정 750개 추출을 재현한다. 비라벨 10,000개는 train 안에서 정답을 보지 않고 무작위 추출하며, 해당 seed에서 n=\allowbreak{}20까지 사용하는 라벨 행과 동일한 특징 벡터를 제외한다. 같은 seed의 비라벨 집합은 n 사이에서 고정한다. test는 일부 표본이 아닌 원래 전체 test이며 seed 43의 라벨과 무관한 순서로 평가한다.

원래 train/\allowbreak{}test 사이의 동일 특징 중복은 유지하고 실행 시 그 수를 기록한다. 실행기는 이 조건을 과거 분할 진단으로 명시해 --allow-test-feature-overlap 옵션을 전달한다. 학습과 비라벨 사이의 동일 특징은 허용하지 않는다. 검토한 기존 분할이므로 독립 평가라고 해석하지 않으며, 정제 효과를 없애기 위해 라벨 충돌 제거도 수행하지 않는다.

\subsection{실제 생성 및 검증}
실제 생성 검증은 두 데이터셋의 n=\allowbreak{}20·seed 43 조건에서 수행했다. 원본 파일을 해시 검증한 뒤 원래 시나리오별 시간 분할을 재생성했으며, 기존 EXP89 클래스 수와 test 행 수를 확인했다. 생성 후 학습 ID를 기존 EXP90의 fixed750\_\allowbreak{}n20\_\allowbreak{}s43과, test ID 및 순서를 EXP89와 원소 단위로 비교하여 모두 일치함을 확인했다. 참라벨/\allowbreak{}비라벨 ID 중복과 동일 특징 벡터 중복도 없었다.

\begingroup
\footnotesize
\begin{longtable}{>{\raggedright\arraybackslash}p{\dimexpr 0.15196\linewidth-2\tabcolsep\relax}>{\raggedleft\arraybackslash}p{\dimexpr 0.15196\linewidth-2\tabcolsep\relax}>{\raggedright\arraybackslash}p{\dimexpr 0.22059\linewidth-2\tabcolsep\relax}>{\raggedleft\arraybackslash}p{\dimexpr 0.15196\linewidth-2\tabcolsep\relax}>{\raggedleft\arraybackslash}p{\dimexpr 0.15196\linewidth-2\tabcolsep\relax}>{\raggedleft\arraybackslash}p{\dimexpr 0.17157\linewidth-2\tabcolsep\relax}}
\toprule
자료 & 정상 라벨 & 공격 라벨 & 비라벨 & 전체 test & NPZ bytes \\ \midrule
\endfirsthead
자료 & 정상 라벨 & 공격 라벨 & 비라벨 & 전체 test & NPZ bytes \\ \midrule
\endhead
{}CIC2018 & 750 & 클래스당 20, 총 120 & 10000 & 4023114 & 157696547 \\
{}ToN-IoT & 750 & 클래스당 20, 총 180 & 10000 & 5504060 & 119953118 \\
\bottomrule
\end{longtable}
\endgroup
정상 클래스 ID는 CIC2018에서 0, ToN-IoT에서 1이며 실행기가 각 JSON에서 읽는다. 특징 열 0과 1의 PROTOCOL·L7\_\allowbreak{}PROTO를 범주형으로 명시한다. 실제 모델의 추론과 성능 평가는 이번 준비 검증에서 수행하지 않았다. 실행 모델의 가중치 갱신·저장 복원 검증은 이전 20261028\_\allowbreak{}exp91 보고서에 구분되어 있다.

단위 검증 4개는 기존 표본 추출 재현, 중첩 예산, 비라벨 특징 겹침 제외, 부족한 pool에서 조용히 예산을 축소하지 않는 동작, 실행기에 정상 ID·범주형 열을 전달하는 경로와 완료 조건의 재사용을 검사했다. 실행기 검증에는 테스트 대역을 사용했으며 성능 근거가 아니다. 실제 NPZ 두 개에 대한 --check도 통과했다. 야간 실행기는 조건별 로그·모델·예측을 기록하고 같은 날 이전 기록을 보존하면서 한국어 보고서를 추가한다.

\subsection{원격 장비 실행 명령}
원격 장비의 프로젝트 루트에서 prepare\_\allowbreak{}inputs.py를 실행하면 입력 경로가 자동으로 정해진다. 환경과 실제 파일 위치만 확인하려면 먼저 --check를 사용한다. 입력을 만든 뒤 등록 실행으로 run\_\allowbreak{}night.py를 실행하며 --model both는 한 GPU에서 1안·2안을 순차 수행한다. 장비를 나누면 --model a와 --model b를 각각 사용한다.


\begin{Verbatim}[breaklines=true,breakanywhere=true,fontsize=\scriptsize]
python scripts/v7/exp91/prepare_inputs.py --check
python scripts/v7/exp91/prepare_inputs.py
python scripts/v7/exp91/run_night.py --check

python scripts/results_versions.py run --version v7 --experiment exp91 \
  --scope tabpfn --label remote_both_s43 \
  --protocol configs/experiments/exp91_remote_protocol.json -- \
  python scripts/v7/exp91/run_night.py --model both --output '{run_dir}'

\end{Verbatim}
GPU를 선택하려면 최상위 명령에 CUDA\_\allowbreak{}VISIBLE\_\allowbreak{}DEVICES를 설정한다. 기본 원본 파일명이 다르면 입력 생성기의 --source에 실제 suite 경로를 지정하며, checkpoint 위치가 다르면 실행기의 --checkpoint를 지정한다. 실행 중에는 등록 경로의 progress.json과 jobs/\allowbreak{}조건명/\allowbreak{}stdout.log를 확인한다. 실패 조건은 오류 로그를 확인한 뒤 같은 입력·설정으로 재개하고, 소스를 변경했다면 새 실행으로 등록한다.

\subsection{근거와 후속 판단}
자료 준비 성공은 분류 이점의 증거가 아니다. 후속 평가는 전체 정상 포함 test의 공격 Macro-F1, 클래스별 오류, 정상 오탐과 전문가 채택률을 확인한다. 두 안은 비라벨 접근 여부가 다르므로 서로의 점수만으로 구조 우위를 판단하지 않고 같은 정보 예산의 기반 모델과 절제를 추가한다. 연구 버전은 v7을 유지하며 논문 원고와 발표자료는 변경하지 않는다.

출처: docs/\allowbreak{}lablog/\allowbreak{}evidence/\allowbreak{}20261029/\allowbreak{}exp91\_\allowbreak{}portable\_\allowbreak{}inputs.json 출처: configs/\allowbreak{}experiments/\allowbreak{}exp91\_\allowbreak{}remote\_\allowbreak{}protocol.json 출처: tabpfn/\allowbreak{}results/\allowbreak{}v7/\allowbreak{}exp91/\allowbreak{}20261029\_\allowbreak{}000418\_\allowbreak{}prepare\_\allowbreak{}portable\_\allowbreak{}inputs\_\allowbreak{}n20\_\allowbreak{}s43\_\allowbreak{}1b05a5c5/\allowbreak{}RUN\_\allowbreak{}META.json


\subsection{원격 누락 파일 오류 정정과 독립 재현 검증}
EXP89의 결과 manifest를 직접 읽던 의존성이 원격 입력 생성을 막았다.
기존 검증은 과거 결과가 있는 작업 폴더에서 수행했으며, Git의 \texttt{data/} 제외 규칙이 결과 경로 안의 data 폴더에도 적용되는 조건을 놓쳤다.
따라서 기존의 원본만으로 실행 가능하다는 설명은 수정 전 코드에는 정확하지 않았다.
수정본은 클래스별 train/test 수와 특징 순서, 원본 해시를 Git에 포함되는 \path{configs/experiments/exp91_input_reference.json}에서 읽는다.
과거 실행 경로는 참조하지 않으며 원본 해시와 분할 검증을 유지한다.

독립 재현 검증은 과거 결과 폴더를 만들지 않은 임시 저장소에 필요한 코드와 새 설정만 복사하고 실제 원본 suite를 연결하여 수행했다.
CIC2018과 ToN-IoT의 n=20, seed 43, 비라벨 10,000개, 전체 test 입력을 다시 생성했으며 두 NPZ 모두 기존 파일과 바이트 단위 해시가 일치했다.
추가 회귀 검증은 과거 결과가 없는 작은 임시 저장소에서 두 데이터셋의 n=2와 n=20 생성 및 완료 입력 재사용을 검사했다.
관련 단위 검증은 총 5개가 통과했으며, 분류 모델의 학습이나 성능 평가는 이번 수정 검증에서 수행하지 않았다.
원격 장비에서는 git pull 후 기존 prepare\_inputs.py 명령을 다시 실행하면 완료된 입력은 재사용하고 미완료 입력을 생성한다.

\sourcepath{docs/lablog/evidence/20261029/exp91_portability_fix.json}
\sourcepath{tabpfn/results/v7/exp91/20261029_003042_portable_no_history_regression_bdcecf7c/validation.json}
\subsection{원격 실행 실패 확인과 오류 노출 개선}
원격 실행의 입력 확인은 통과했지만 사용자 제공 상위 로그에 나타난 세 실행은 각각 10개 조건 모두 실패했다.
장비의 확인 시각은 2026-10-29 00:58 KST이며 입력 10개와 checkpoint 경로는 존재했다.
\texttt{--check}는 파일과 메타데이터만 검사하므로 GPU 및 TabPFN 실행 성공을 뜻하지 않는다.
추가로 제공된 최신 실행의 CIC2018 n=2 작업 로그에서 GPU 초기화가 드라이버 호환 오류로 실패했음을 확인했다.
로그의 드라이버 API 버전은 12040이며 실행 Python은 다른 프로젝트의 Python 3.14 가상환경이었다.
보고서 생성은 별도로 \texttt{lxml} 패키지 누락 때문에 실패했다.

실행기는 첫 실패에서 중단하고 개별 로그 마지막 80줄과 보고서 생성 예외를 상위 로그에 출력하도록 개선했다.
\texttt{--keep-going}을 명시하면 실패 후 다른 조건을 계속 수행할 수 있으며, 기본 동작은 공통 오류의 반복 실행을 막는다.
기존 입력은 재사용하고 공용 설치기의 Python 3.12·드라이버별 CUDA 빌드 선택으로 실행 환경을 교정한다.
CUDA 12.4 지원 드라이버의 선택 결과는 PyTorch 2.6.0+cu124이며 공용 requirements에 \texttt{lxml}과 \texttt{markdown-it-py}를 추가했다.
\texttt{--check}와 실제 실행 직전 검사에 패키지 import 및 장치 행렬 연산을 추가해 데이터 전체 로딩 전에 오류를 발견한다.
이 검사는 모델 fit이나 성능 검증을 대신하지 않는다.

로컬 단위 검증 17개가 통과했고, 추가한 검증은 실패 즉시 중단·명시적 계속 실행·오류 내용 노출·완료 표식 미생성 및 패키지 누락과 드라이버 오류 동시 보고를 확인했다.
빈 결과의 보고서 생성은 격리된 로컬 경로에서 성공했다.
원격 장비의 수정 환경 설치와 모델 실행 완료는 이 로컬 검증에 포함되지 않으며 재실행 후 확인해야 한다.

\sourcepath{docs/lablog/evidence/20261029/exp91_remote_failure_diagnostics.json}

% EXP91_RUN: 20261028_161815_remote_both_s43_fixedenv_ecd34e66
\subsection{실행 20261028\_\allowbreak{}161815\_\allowbreak{}remote\_\allowbreak{}both\_\allowbreak{}s43\_\allowbreak{}fixedenv\_\allowbreak{}ecd34e66}
두 제안 분류기의 실제 자료 실행을 기록한다. 완료 조건은 10/\allowbreak{}10개이며 실행 모델은 both, seed는 [43], 공격별 라벨 예산은 [2, 3, 5, 10, 20]이다. 정상 라벨은 750개, 비라벨 학습 행은 10000개다. 실제 실행 경로는 tabpfn/\allowbreak{}results/\allowbreak{}v7/\allowbreak{}exp91/\allowbreak{}20261028\_\allowbreak{}161815\_\allowbreak{}remote\_\allowbreak{}both\_\allowbreak{}s43\_\allowbreak{}fixedenv\_\allowbreak{}ecd34e66이다.

원래 시나리오별 시간 분할과 정제 전 특징을 사용한다. 각 모델은 정상 포함 전체 원래 test를 평가한다. 기존에 검토한 분할이며 train/\allowbreak{}비라벨과 test 사이의 동일 특징 중복을 유지하고 그 수를 감사 파일에 기록한다. 따라서 이 결과를 독립 분할 검증이나 문헌 대비 우위로 해석하지 않는다.

\begingroup
\footnotesize
\begin{longtable}{>{\raggedright\arraybackslash}p{\dimexpr 0.32384\linewidth-2\tabcolsep\relax}>{\raggedright\arraybackslash}p{\dimexpr 0.12764\linewidth-2\tabcolsep\relax}>{\raggedleft\arraybackslash}p{\dimexpr 0.15295\linewidth-2\tabcolsep\relax}>{\raggedleft\arraybackslash}p{\dimexpr 0.13397\linewidth-2\tabcolsep\relax}>{\raggedleft\arraybackslash}p{\dimexpr 0.12764\linewidth-2\tabcolsep\relax}>{\raggedleft\arraybackslash}p{\dimexpr 0.13397\linewidth-2\tabcolsep\relax}}
\toprule
조건 & 모델 & 공격 Macro-F1 & 정상 오탐률 & 학습 초 & 추론 초 \\ \midrule
\endfirsthead
조건 & 모델 & 공격 Macro-F1 & 정상 오탐률 & 학습 초 & 추론 초 \\ \midrule
\endhead
{}cic2018\_\allowbreak{}original\_\allowbreak{}fixed750\_\allowbreak{}n2\_\allowbreak{}s43\_\allowbreak{}u10000\_\allowbreak{}full & a & 0.442659 & 0.005908 & 32.642 & 5.200 \\
{}cic2018\_\allowbreak{}original\_\allowbreak{}fixed750\_\allowbreak{}n2\_\allowbreak{}s43\_\allowbreak{}u10000\_\allowbreak{}full & b & 0.412322 & 0.000691 & 0.867 & 1087.175 \\
{}toniot\_\allowbreak{}original\_\allowbreak{}fixed750\_\allowbreak{}n2\_\allowbreak{}s43\_\allowbreak{}u10000\_\allowbreak{}full & a & 0.246397 & 0.024142 & 71.098 & 10.729 \\
{}toniot\_\allowbreak{}original\_\allowbreak{}fixed750\_\allowbreak{}n2\_\allowbreak{}s43\_\allowbreak{}u10000\_\allowbreak{}full & b & 0.168799 & 0.001189 & 1.118 & 929.758 \\
{}cic2018\_\allowbreak{}original\_\allowbreak{}fixed750\_\allowbreak{}n3\_\allowbreak{}s43\_\allowbreak{}u10000\_\allowbreak{}full & a & 0.610190 & 0.009326 & 32.869 & 5.122 \\
{}cic2018\_\allowbreak{}original\_\allowbreak{}fixed750\_\allowbreak{}n3\_\allowbreak{}s43\_\allowbreak{}u10000\_\allowbreak{}full & b & 0.554336 & 0.000662 & 12.880 & 596.751 \\
{}toniot\_\allowbreak{}original\_\allowbreak{}fixed750\_\allowbreak{}n3\_\allowbreak{}s43\_\allowbreak{}u10000\_\allowbreak{}full & a & 0.339983 & 0.029434 & 32.555 & 7.447 \\
{}toniot\_\allowbreak{}original\_\allowbreak{}fixed750\_\allowbreak{}n3\_\allowbreak{}s43\_\allowbreak{}u10000\_\allowbreak{}full & b & 0.368991 & 0.023559 & 25.041 & 830.413 \\
{}cic2018\_\allowbreak{}original\_\allowbreak{}fixed750\_\allowbreak{}n5\_\allowbreak{}s43\_\allowbreak{}u10000\_\allowbreak{}full & a & 0.662905 & 0.007442 & 32.364 & 5.258 \\
{}cic2018\_\allowbreak{}original\_\allowbreak{}fixed750\_\allowbreak{}n5\_\allowbreak{}s43\_\allowbreak{}u10000\_\allowbreak{}full & b & 0.675158 & 0.000706 & 14.113 & 607.137 \\
{}toniot\_\allowbreak{}original\_\allowbreak{}fixed750\_\allowbreak{}n5\_\allowbreak{}s43\_\allowbreak{}u10000\_\allowbreak{}full & a & 0.427048 & 0.051074 & 32.139 & 7.520 \\
{}toniot\_\allowbreak{}original\_\allowbreak{}fixed750\_\allowbreak{}n5\_\allowbreak{}s43\_\allowbreak{}u10000\_\allowbreak{}full & b & 0.425543 & 0.043738 & 45.454 & 829.157 \\
{}cic2018\_\allowbreak{}original\_\allowbreak{}fixed750\_\allowbreak{}n10\_\allowbreak{}s43\_\allowbreak{}u10000\_\allowbreak{}full & a & 0.651179 & 0.011648 & 32.506 & 5.130 \\
{}cic2018\_\allowbreak{}original\_\allowbreak{}fixed750\_\allowbreak{}n10\_\allowbreak{}s43\_\allowbreak{}u10000\_\allowbreak{}full & b & 0.496802 & 0.000787 & 32.876 & 1197.744 \\
{}toniot\_\allowbreak{}original\_\allowbreak{}fixed750\_\allowbreak{}n10\_\allowbreak{}s43\_\allowbreak{}u10000\_\allowbreak{}full & a & 0.497348 & 0.085945 & 32.171 & 7.418 \\
{}toniot\_\allowbreak{}original\_\allowbreak{}fixed750\_\allowbreak{}n10\_\allowbreak{}s43\_\allowbreak{}u10000\_\allowbreak{}full & b & 0.451340 & 0.072423 & 60.489 & 1655.424 \\
{}cic2018\_\allowbreak{}original\_\allowbreak{}fixed750\_\allowbreak{}n20\_\allowbreak{}s43\_\allowbreak{}u10000\_\allowbreak{}full & a & 0.632698 & 0.041435 & 32.679 & 5.056 \\
{}cic2018\_\allowbreak{}original\_\allowbreak{}fixed750\_\allowbreak{}n20\_\allowbreak{}s43\_\allowbreak{}u10000\_\allowbreak{}full & b & 0.695031 & 0.000724 & 31.710 & 618.344 \\
{}toniot\_\allowbreak{}original\_\allowbreak{}fixed750\_\allowbreak{}n20\_\allowbreak{}s43\_\allowbreak{}u10000\_\allowbreak{}full & a & 0.489762 & 0.099168 & 31.817 & 7.444 \\
{}toniot\_\allowbreak{}original\_\allowbreak{}fixed750\_\allowbreak{}n20\_\allowbreak{}s43\_\allowbreak{}u10000\_\allowbreak{}full & b & 0.492375 & 0.082536 & 87.111 & 2517.011 \\
\bottomrule
\end{longtable}
\endgroup
실패 상태는 []이며 각 조건의 stdout.log에서 원인을 확인한다. 실패 조건은 동일한 설정으로 재개하고, 설정이나 소스를 수정할 때에는 새 실행으로 등록한다. 후속 분석은 클래스별 오류와 전문가 채택 여부, 정상 오탐, 같은 정보 예산의 기반 모델 및 절제 대조를 확인한다.

출처: tabpfn/\allowbreak{}results/\allowbreak{}v7/\allowbreak{}exp91/\allowbreak{}20261028\_\allowbreak{}161815\_\allowbreak{}remote\_\allowbreak{}both\_\allowbreak{}s43\_\allowbreak{}fixedenv\_\allowbreak{}ecd34e66/\allowbreak{}summary.json 및 각 jobs/\allowbreak{} 하위 config.json·input\_\allowbreak{}manifest.json·감사 파일.


\end{document}
